\chapter{Conclusion}\label{ch:conclusion}
Nous voulions savoir ce que l'entraînement par renforcement apprend à un agent de recherche. Comparé à des harnais fixes, l'agent Search-R1 répond mieux ; comparé à son propre modèle de départ, il répond mieux au premier essai sur les questions qui obligent à enchaîner les recherches, mais il ne résout pas davantage de questions quand on lui laisse plusieurs essais. Le renforcement l'a rendu plus fiable, sans élargir ce qu'il sait faire. Ce gain tient d'abord à la recherche : l'agent trouve plus souvent les deux paragraphes nécessaires. Il lit aussi mieux, mais renonce trop souvent à répondre pour que cela se voie. Ces résultats corrigent notre première conclusion, qui attribuait l'avantage à la lecture à partir de modèles de départ différents. Ils déplacent la question pratique : si le renforcement fiabilise surtout la recherche, il reste à savoir si une voie moins coûteuse, l'apprentissage supervisé sur les bonnes traces, en obtient autant.
